\subsection{Materiales}

Para la realización del experimento se utilizó un bulbo de vacío tipo 1V2 montado en un módulo de conexión, una fuente de corriente directa variable de \SIrange{0}{5}{V} y \SI{2}{A} para alimentar el filamento, una fuente de corriente directa variable de \SIrange{0}{60}{V} y \SI{1}{A} para polarizar el ánodo, y un multímetro configurado como miliamperímetro para medir la corriente de ánodo. El voltaje del filamento se mantuvo fijo en \SI{0.63}{V}, mientras que el voltaje de ánodo se varió de forma controlada durante la adquisición de datos.

\subsection{Diseño experimental}

El montaje consistió en polarizar el filamento del bulbo 1V2 para producir emisión termoiónica y aplicar una diferencia de potencial positiva entre el ánodo y el cátodo para colectar los electrones emitidos. La corriente de ánodo $I_a$ se midió conectando el multímetro en serie con el ánodo. La fuente del filamento se utilizó únicamente para calentar el cátodo, por lo que su voltaje se mantuvo constante durante toda la toma de datos. La fuente del ánodo proporcionó la variable independiente del experimento, $V_a$, mientras que la corriente $I_a$ fue la variable dependiente.

Antes de iniciar las mediciones se verificó que la fuente de ánodo estuviera en \SI{0}{V} y que la fuente del filamento estuviera ajustada al valor indicado. Después de conectar las fuentes y el multímetro al módulo del bulbo, se dejó calentar el filamento durante aproximadamente cinco minutos para estabilizar la emisión termoiónica.

\subsection{Procedimiento}

Primero se conectó el bulbo 1V2 al módulo experimental siguiendo el diagrama de polarización del manual de la práctica. El filamento se alimentó con un voltaje constante de \SI{0.63}{V}, usando la fuente de bajo voltaje. La fuente de ánodo se conectó entre el ánodo y el cátodo, cuidando la polaridad indicada para que el ánodo quedara a potencial positivo respecto al filamento. El multímetro se conectó en la escala de corriente adecuada para registrar la corriente de ánodo $I_a$ en miliamperes.

Una vez estabilizado el filamento, se varió el voltaje de ánodo desde \SI{0}{V} hasta \SI{20}{V} en incrementos de \SI{0.5}{V}. Para cada valor de $V_a$ se registró la corriente correspondiente $I_a$. Posteriormente, el voltaje de ánodo se continuó variando desde \SI{20}{V} hasta \SI{40}{V} en incrementos de \SI{1}{V}, registrando nuevamente la corriente de ánodo para cada punto. Durante toda la medición se mantuvo constante el voltaje de filamento.

